\documentclass{standalone}
%\documentclass[11pt,hyperref={colorlinks=false,runcolor=blue}]{beamer}

%\usepackage{tgheros}
%\renewcommand*\familydefault{\sfdefault} %% Only if the base font of the document is to be sans serif
%\usepackage[T1]{fontenc}

\usepackage{arev}
\usepackage[T1]{fontenc}

%\usepackage{PTSansCaption} 
%\renewcommand*\familydefault{\sfdefault} %% Only if the base font of the document is to be sans serif
%\usepackage[T1]{fontenc}


\usepackage{tikz}
\usetikzlibrary{backgrounds,fit,calc,shadows,chains,ocgx,shapes.geometric,arrows,positioning}
%\usepackage{microtype}
\usepackage{url}
%\usepackage{fancyvrb}
%\usepackage{listings}
\usepackage{pgfplots}


\definecolor{mygreen}{RGB}{127,196,79}
\definecolor{myblue}{RGB}{72,175,194}
\definecolor{mydarkblue}{RGB}{45,108,134}
\definecolor{mylightblue}{RGB}{181,223,230}
\definecolor{mybrown}{RGB}{178,140,110}
\definecolor{mypink}{RGB}{173,110,207}
\definecolor{myorange}{RGB}{238,120,0}
\definecolor{myred}{RGB}{186,62,93}
\definecolor{mylandfill}{RGB}{164,141,87}
\definecolor{mysludge}{RGB}{210,190,135}
\definecolor{myocean}{RGB}{70,150,155}

\definecolor{lightblue}{RGB}{212,226,237}
\definecolor{lightbrown}{RGB}{247,240,220}
\definecolor{lightgreen}{RGB}{232,247,240}
\definecolor{lightgray}{RGB}{237,237,237}

\definecolor{fillgreen}{RGB}{164,218,143}
\definecolor{fillyellow}{RGB}{255,237,171}
\definecolor{fillblue}{RGB}{167,223,226}
\definecolor{fillpurple}{RGB}{226,197,236}
\definecolor{fillorange}{RGB}{251,216,181}
\definecolor{fillpurple}{RGB}{231,208,239}
\definecolor{fillstorm}{RGB}{173,203,228}
\definecolor{filllandfill}{RGB}{234,215,164}
\definecolor{fillsludge}{RGB}{238,220,179}
\definecolor{fillocean}{RGB}{103,181,183}

% USE readarray  TO GET THE DATA INTO A \def
\usepackage{readarray}
\readarraysepchar{;}

\usepackage{etoolbox}

% for too large figure margins because of bounding boxes around the tikz figure
\usepackage{adjustbox}

\begin{document}

% defining a new command for plotting the barplots
% needs to be used inside a tikzpicture environment
% #1 first variable is the name of the flow to look up in the file
% #2 is the name of the button that will switch it on
\newcommand{\barplot}[4]{
\begin{scope}[xshift=#3 cm, yshift=#4 cm,ocg={name=#2,ref=#2,status=invisible}]

  %\hideocg{s0a s0b s1 s2 s3 s4 s5 s6 s7 s8 s9 s10 s11 s12 s13 s14 s15 s16 s17 s18 s19 s20 s21}{hide}  
  
  \draw[draw=black,fill=white,drop shadow={fill=black}] (-2,-1) rectangle ++(11.5,9);
  
  \begin{axis}[
    title=\textbf{#1} ,
    width=10cm,
    ylabel={Flow (t)},
    %enlargelimits=0.1,
    ymin=0,
    bar width=7mm,
    grid=major,
    symbolic x coords={LDPE,HDPE,PP,PS,EPS,PVC,PET}
    ]
    
  \addplot[ybar,fill=red!70] 
    coordinates {(LDPE,\LDPEmean{#1})};
    
  \addplot[ybar,fill=orange!72] 
    coordinates {(HDPE,\HDPEmean{#1})};
    
  \addplot[ybar,fill=yellow!75] 
    coordinates {(PP,\PPmean{#1})};
    
  \addplot[ybar,fill=lime!75] 
    coordinates {(PS,\PSmean{#1})};
    
  \addplot[ybar,fill=cyan!72] 
    coordinates {(EPS,\EPSmean{#1})};
    
  \addplot[ybar,fill=blue!70] 
    coordinates {(PVC,\PVCmean{#1})};
    
  \addplot[ybar,fill=violet!70] 
    coordinates {(PET,\PETmean{#1})};

  \addplot[draw=black,ultra thick] coordinates {(LDPE,\LDPElow{#1}) (LDPE,\LDPEhigh{#1})};
  \addplot[draw=black,ultra thick] coordinates {(HDPE,\HDPElow{#1}) (HDPE,\HDPEhigh{#1})};
  \addplot[draw=black,ultra thick] coordinates {(PP,\PPlow{#1}) (PP,\PPhigh{#1})};
  \addplot[draw=black,ultra thick] coordinates {(PS,\PSlow{#1}) (PS,\PShigh{#1})};
  \addplot[draw=black,ultra thick] coordinates {(EPS,\EPSlow{#1}) (EPS,\EPShigh{#1})};
  \addplot[draw=black,ultra thick] coordinates {(PVC,\PVClow{#1}) (PVC,\PVChigh{#1})};
  \addplot[draw=black,ultra thick] coordinates {(PET,\PETlow{#1}) (PET,\PEThigh{#1})};

  \end{axis}  
  
  \node at (4.25,6.5) {\textbf{\tottext{#1} t}};
  
\end{scope}
}

% create interactive button function
\tikzset{%
  button on/.style={%
    draw,circle,drop shadow,
    line width=1pt,
    fill=orange!50!red!50,
    switch ocg with mark on={#1}{},
  },
}